\section{The \texttt{Fluorescence} McXtrace Component}
Sample model handling absorption, fluorescence, Compton and Rayleigh scattering.

\subsection*{Identification}
\begin{itemize}
  \item \textbf{Author:} E. Farhi
  \item \textbf{Origin:} Synchrotron SOLEIL
  \item \textbf{Date:} March 2022
\end{itemize}

\subsection*{Description}
Sample that models many photon-matter interactions:

\begin{verbatim}
- absorption          (photon excites an electron and creates a hole)
- fluorescence        (excited electrons emit light while falling into lower states)
- Compton scattering  (inelastic, transfer some energy to an electron,    incoherent)
- Rayleigh scattering (elastic,   dipolar radiation of excitated electrons, coherent)
\end{verbatim}

Use the \textbf{FluoPowder} sample component to properly handle powder diffraction with fluorescence.

The 'material' specification is given as a chemical formulae, e.g. "LaB6". It may also be given as a file name (CIF/LAU/LAZ/FullProf format) in which case the formulae is guessed (but may be approximative). Atoms from Z=5 to Z=90 are handled.

By setting the 'order' to 1, the absorption along the scattered path is handled. A higher 'order' will handle multiple scattering events, and final absorption. For instance, a value order\textgreater{}=2 handles e.g. fluorescence iterative cascades in the material. Leaving 'order=0' handles the single scattering only.

Example: Fluorescence(material="LaB6", xwidth=0.001,yheight=0.001,zdepth=0.0001, p\_interact=0.99, target\_index=1, focus\_xw=0.0005, focus\_yh=0.0005)

\textbf{Sample shape:} Sample shape may be a cylinder, a sphere, a box or any other shape

\begin{verbatim}
box/plate:       xwidth x yheight x zdepth (thickness=0)
\end{verbatim}

hollow box/plate:xwidth x yheight x zdepth and thickness\textgreater{}0

\begin{verbatim}
cylinder:        radius x yheight (thickness=0)
\end{verbatim}

hollow cylinder: radius x yheight and thickness\textgreater{}0

\begin{verbatim}
sphere:          radius (yheight=0 thickness=0)
hollow sphere:   radius and thickness>0 (yheight=0)
any shape:       geometry=OFF file
\end{verbatim}

The complex geometry option handles any closed non-convex polyhedra. It computes the intersection points of the photon ray with the object transparently, so that it can be used like a regular sample object. It supports the OFF, PLY and NOFF file format but not COFF (colored faces). Such files may be generated from XYZ data using: qhull \textless{} coordinates.xyz Qx Qv Tv o \textgreater{} geomview.off or powercrust coordinates.xyz and viewed with geomview or java -jar jroff.jar (see below). The default size of the object depends of the OFF file data, but its bounding box may be resized using xwidth,yheight and zdepth.

\textbf{Concentric components:} This component has the ability to contain other components when used in hollow cylinder geometry (namely sample environment, e.g. cryostat and furnace structure). Such component 'shells' should be split into input and output side surrounding the 'inside' components. First part must then use 'concentric=1' flag to enter the inside part. The component itself must be repeated to mark the end of the concentric zone. The number of concentric shells and number of components inside is not limited.

COMPONENT F\_in = Fluorescence(material="Al", concentric=1, ...) AT (0,0,0) RELATIVE sample\_position

COMPONENT something\_inside ... // e.g. the sample itself or other materials

COMPONENT F\_out = COPY(F\_in)(concentric=0) AT (0,0,0) RELATIVE sample\_position

\textbf{Enhancing computation efficiency:} * An important option to enhance statistics is to set 'p\_interact' to, say, 30 percent (0.3) in order to force a fraction of the beam to scatter. This will result on a larger number of scattered events, retaining intensity.

In addition, it may be desirable to define a 'target' for the fluorescence processes via e.g. the 'target\_index' and the 'focus\_xw / focus\_yh' options. This target should e.g. be the SDD area.

This sample component can advantageously benefit from the SPLIT feature, e.g. SPLIT COMPONENT sample = Fluorescence(...)

The computation is made via the XRayLib (apt install libxrl-dev) https://github.com/tschoonj/xraylib.

\subsection*{Input parameters}
Parameters in \textbf{boldface} are required; the others are optional.

\begin{longtable}{p{0.22\textwidth}p{0.12\textwidth}p{0.46\textwidth}p{0.14\textwidth}}
\toprule
\textbf{Name} & \textbf{Unit} & \textbf{Description} & \textbf{Default} \\
\midrule
\endhead
geometry & str & Name of an Object File Format (OFF) or PLY file for complex geometry. The OFF/PLY file may be generated from XYZ coordinates using qhull/powercrust. & 0 \\
radius & m & Outer radius of sample in (x,z) plane. cylinder/sphere. & 0 \\
thickness & m & Thickness of hollow sample Negative value extends the hollow volume outside of the box/cylinder. & 0 \\
xwidth & m & Width for a box sample shape. & 0 \\
yheight & m & Height of sample in vertical direction for box/cylinder shapes. & 0 \\
zdepth & m & Depth for a box sample shape. & 0 \\
concentric & 1 & Indicate that this component has a hollow geometry and may contain other components. It should then be duplicated after the inside part (only for box, cylinder, sphere). & 0 \\
material & str & Chemical formulae, e.g. "LaB6", "Pb2SnO4". If may also be a CIF/LAZ/LAU file. & "LaB6" \\
packing\_factor & 1 & How dense is the material compared to bulk 0-1. & 0 \\
rho & \AA{}-3 & Density of scattering elements (nb atoms/unit cell V\_0). & 0 \\
density & g/cm$^{3}$ & Density of material. V\_rho=density/weight/1e24*N\_A. & 0 \\
weight & g/mol & Atomic/molecular weight of material. & 0 \\
p\_interact & 1 & Force a given fraction of the beam to scatter, keeping intensity right, to enhance small signals (-1 inactivate). & 0 \\
target\_x & m & Position of target to focus at, along X. & 0 \\
target\_y & m & Position of target to focus at, along Y. & 0 \\
target\_z & m & Position of target to focus at, along Z. & 0 \\
focus\_r & m & Radius of disk containing target. Use 0 for full space. & 0 \\
focus\_xw & m & Horiz. dimension of a rectangular area. & 0 \\
focus\_yh & m & Vert.  dimension of a rectangular area. & 0 \\
focus\_aw & deg & Horiz. angular dimension of a rectangular area. & 0 \\
focus\_ah & deg & Vert.  angular dimension of a rectangular area. & 0 \\
target\_index & 1 & Relative index of component to focus at, e.g. next is +1. & 0 \\
flag\_compton & 1 & When 0, the Compton  scattering is ignored. & 1 \\
flag\_rayleigh & 1 & When 0, the Rayleigh scattering is ignored. & 1 \\
flag\_lorentzian & 1 & When 1, the line shapes are assumed to be Lorentzian, else Gaussian. & 0 \\
flag\_kissel & 1 & When 1 (slower), handle M-lines XRF from Kissel for Z\textgreater{}=52 Te (else only K and L-lines). & 0 \\
flag\_low\_z & 1 & When 1, enhances low concentration atoms, else uses cross-sections weighting. & 0 \\
order & 1 & Limit multiple fluorescence up to given order. Last iteration is absorption only. & 1 \\
\bottomrule
\end{longtable}

\subsection*{Links}
\begin{itemize}
  \item Component source code found in file \texttt{Fluorescence.comp}.
  \item The XRayLib https://github.com/tschoonj/xraylib http://dx.doi.org/10.1016/j.sab.2011.09.011
  \item Fluorescence https://en.wikipedia.org/wiki/Fluorescence
  \item Rayleigh https://en.wikipedia.org/wiki/Rayleigh\_scattering
  \item Compton https://en.wikipedia.org/wiki/Compton\_scattering
  \item X-ray absorption edges http://skuld.bmsc.washington.edu/scatter/AS\_periodic.html
  \item X-ray fluorescence spectra http://www.xrfresearch.com/xrf-spectra/
  \item X-ray edges and fluo lines https://physics.nist.gov/PhysRefData/XrayTrans/Html/search.html
\end{itemize}
\IfFileExists{samples/Fluorescence_static.tex}{\input{samples/Fluorescence_static.tex}}{}